```latex
\documentclass[12pt]{book}
\usepackage[utf8]{inputenc}
\usepackage[T1]{fontenc}
\usepackage{lmodern}
\usepackage{geometry}
\usepackage{xcolor}
\usepackage{titlesec}
\usepackage{setspace}
\usepackage{fancyhdr}
\usepackage{mdframed}
\usepackage{needspace}
\usepackage{enumitem}
\usepackage{hyperref}
\usepackage{array}
\usepackage{booktabs}
\usepackage{tabularx}
\usepackage{longtable}
\usepackage{calc}
\usepackage{footnote}
\makesavenoteenv{tabular}
\makesavenoteenv{longtable}
\usepackage{pdflscape}
\usepackage{multicol}
\usepackage{graphicx}
\hypersetup{
    colorlinks=true,
    linkcolor=darkblue,
    urlcolor=darkblue,
    pdftitle={The Eastern Realms},
    pdfauthor={Dusana of the Raven Road}
}

% ------------------- Custom commands -------------------
\newcommand{\tag}[1]{\textsc{#1}}
\newcommand{\clk}[1]{\textit{[#1]}}
\newcommand{\stringitem}[1]{\textit{#1}}
\newcommand{\dusana}[1]{%
    \begin{quote}\itshape #1 —~Dusana of the Raven Road\end{quote}}
\newcommand{\letter}[2]{%
    \begin{quote}\small\itshape #1 \hfill --- #2\end{quote}}

% ------------------- Styling -------------------
\titleformat{\chapter}{\Huge\bfseries}{\thechapter}{20pt}{}
\titleformat{\section}{\Large\bfseries}{\thesection}{15pt}{}
\setlength{\parindent}{0pt}
\setlength{\parskip}{0.8ex}

% ------------------- Infobox styling -------------------
\newmdenv[
    linecolor=black,
    backgroundcolor=gray!5,
    innertopmargin=8pt,
    innerbottommargin=8pt,
    innerrightmargin=8pt,
    innerleftmargin=8pt,
    skipabove=6pt,
    skipbelow=6pt,
    font=\small
]{infobox}

% ------------------- Document -------------------
\begin{document}

\title{
    \Huge\textbf{The Eastern Realms}\\[0.5em]
    \Large Lotus, Blade, and Bell
}
\author{
    Set down by Dusana of the Raven Road\\
    \small Under the lamp, bread and salt
}
\date{
    \small The ink is cut with monsoon rain
}
\maketitle

\tableofcontents

\chapter{Preface: The East That Is Not One Thing}

The maps of the empire draw the East as a single sweep—Sihai, Nihori, Ayokha—all pressed together like tiles on a game board. The roads know better. Sihai measures time in dynasties; Nihori counts it in storms. Ayokha negotiates between them with the patience of the monsoon. And beneath them all, the older peoples keep their own thresholds, their own debts, their own ways of not being seen.

This expansion presents the Eastern Realms not as a backdrop but as a web of obligations, courtesies, and carefully maintained boundaries. You will find travel decks for Sihai, Nihori, and Ayokha; Patrons who speak in ink and ash; and—most importantly—the survival tactics of the small peoples who have lived here longest.

\vspace{1em}
\dusana{I spent a season among the silk farms of Ayokha, and another among the bell‑towers of the Inland Sea. The Lethai there do not speak of the Valewood; they speak of a different Fall, one that left them custodians of a forgotten tide between mortals and the fair folk. The gnomes do not hide from the fae—they live with them, borrowing sugar and never returning it. And the halflings? They still pour the ninth cup and hope the Hollow is not hungry. The road changes the walker. Walk carefully. Count to eight. Light a lamp before dusk.}

\section{How to Use This Book}

Each of the three major realms—Sihai, Nihori, Ayokha—receives a travel deck (Spade=place, Heart=actor, Club=pressure, Diamond=leverage) following the standard \textit{Fate’s Edge} travel procedure. The highest rank sets clock size (4/6/8/10). The Inland Sea provides a shared deck for crossings.

New cultural mechanics appear for the Eastern threshold peoples: the \textbf{Lethai‑thora} (spirit‑bound memory‑keepers), \textbf{Lethai‑al} (storm‑touched wanderers), \textbf{Lethai‑ar} (river‑bound oath‑keepers), \textbf{Aelinnel} (bell‑tending gnomes), and \textbf{Aelaerem} (halflings who court the Hollow). Each includes etiquette hooks, strings, and a unique survival tactic.

Finally, you will find eight new Patrons—spirits of grain, sun, storm, rivers, commerce, war, foxes, and the sky—each with Rites and obligations suited to the Eastern threshold.

\chapter{The Threshold Peoples of the East}

The Eastern Realms are not only the domain of human empires. Beneath the Mandate, the storm, and the monsoon live those who remember older pacts. They are the descendants of spirits who chose mortality, the keepers of thresholds, and the neighbors of the Hollow. Each has a unique way of surviving—and thriving—between empires.

\section{Lethai‑thora — Spirit‑Bound Memory‑Keepers of Sihai}

\noindent \textit{We remember the tide that never rose, the song that never sang.}

Long before the first human emperor, the Lethai‑thora served a celestial court of fey beings—\textit{xian} in the old tongue, \textit{kami} in the northern isles. Their blood still carries the echo of that older world. When the old gods fell silent, the Lethai‑thora fell with them, becoming the bridge between mortals and the fair folk. They dwell in hidden pavilions, preserving “forgotten ledgers” that record what the Mandate of Heaven prefers to forget.

Unlike their kin in the western Valewood, the Lethai‑thora lean toward the \textbf{Gift of the Mind}. They are keepers of context, scholars of lost ritual, and the only ones who can still read the script of the dynasty that never was. That knowledge is a burden: each secret they preserve slowly erases a piece of themselves. An Lethai‑thora who guards a forbidden truth may find that she can no longer recall the color of her mother’s dress—a small price, she says, for keeping the ledger whole.

\vspace{0.5em}
\dusana{A Lethai‑thora once let me read a scroll that had no words. “It’s not blank,” she said. “The words are waiting for the right season.” She refused to tell me when the season would come.}

\letter{“We do not burn the dangerous scrolls. We bury them in plain sight, on shelves where no one thinks to look. The Emperor’s Censor once stood in our archive and asked for a tax ledger. We gave him one. He did not notice the shelf behind him held the true history of his own grandfather’s death. Silence is not emptiness. It is a door that only opens inward.”}{— Lian, Keeper of the Seventh Gutter, from a letter found pinned to a corkboard in a tea house}

\begin{infobox}
\textbf{Cultural Mechanics}
\begin{itemize}
    \item \textbf{Silence as Sanctuary:} Once per scene, when asked a direct question, you may answer with a silence. That silence counts as a successful Deception roll against any attempt to extract the truth—but the GM gains 1 Story Beat (the silence has weight).
    \item \textbf{Context Keys:} To read an old text or recall a forbidden history, you need three keys: place, season, and witness. Without all three, roll at disadvantage. A Lethai‑thora may create a temporary Context Key by spending 1 Boon and describing a small ritual (lighting incense, touching a stone, speaking a name).
    \item \textbf{The Keeper’s Toll:} When you use a Lethai‑thora Rite or consult a hidden archive, mark 1 Fatigue. The knowledge costs something—a moment of warmth, a pleasant memory, the scent of rain.
\end{itemize}
\textbf{Etiquette Hooks}
\begin{itemize}
    \item Offer ink before speech – \texttt{Position +1} in Lethai‑thora halls.
    \item Never touch a scroll without announcing your lineage (even a lie improves DV by 1).
    \item A bow held three seconds longer than expected means “I will remember this slight.”
    \item At dusk, pour a drop of water onto the ground before speaking of the dead.
\end{itemize}
\textbf{Strings:} Forbidden scroll (one truth no one wants spoken), ink‑stone (right to consult a hidden archive), silence‑token (a debt of secrecy owed to you), gloss‑paper (turns one forgotten fact into a clue).
\end{infobox}

\section{Lethai‑al — Storm‑Touched Wanderers of Nihori}

\noindent \textit{Our empire fell to the sea. We learned to read storms instead of dynasties.}

The Lethai‑al of Nihori once guarded a drowned kingdom—an archipelago that sank when a volcanic name was spoken. Their blood still carries the salt of that ancient catastrophe. Now they are wanderers in the oldest sense: masterless, bound to no lord, but fiercely bound to each other and to the storm spirits that swirl above the waves. They dwell in hidden mountain groves, emerging to duel winds, pacify storm‑spirits, and occasionally lift burdens from over‑taxed merchants.

They lean toward the \textbf{Gift of the Body}, but their strength is endurance rather than brute force. A Lethai‑al can walk three days without sleep, hold breath beneath a tidal wave, and survive a fall that would shatter mortal bones. They say this is the blessing of their drowned ancestors, who learned to live where the sea did not want them.

\vspace{0.5em}
\dusana{“A Lethai‑al once told me that the secret to reading a storm is to stop reading and start listening. ‘The wind has a voice,’ she said. ‘You just have to forget your name long enough to hear it.’”}

\letter{“My grandmother used to say: ‘The sea does not forgive, but it does forget. If you can wait long enough, the waves will erase your debt.’ She waited eighty years for a storm that never came. I do not wait. I walk into the rain and let it decide.”}{— Kenji, who left his sword at his mother’s grave and never drew another}

\begin{infobox}
\textbf{Cultural Mechanics}
\begin{itemize}
    \item \textbf{Storm‑Reading:} Once per journey, you may predict the weather for the next 24 hours with perfect accuracy (Wits + Survival vs. DV 3). On success, choose one leg to ignore weather penalties. On a miss, the storm finds you anyway, and the GM gains 1 Story Beat to spend on a thunderclap or flash flood.
    \item \textbf{Broken Oath Compensation:} When you break a promise (even a minor one), you must immediately offer compensation—a gift, a service, or a truth. If you refuse, you gain the condition \textit{Ronin’s Shame} (–1 die to all social rolls until you make amends). A Lethai‑al who accepts a geas may gain +1 die to actions that serve it, but breaking it costs 2 Harm (Spirit).
    \item \textbf{Tidal Resilience:} Once per scene, when you would take Harm from a water‑based or storm‑based effect, you may reduce it by 1 level (to a minimum of 0). Mark 1 Fatigue.
\end{itemize}
\textbf{Etiquette Hooks}
\begin{itemize}
    \item Leave a single arrow behind when you leave a camp (\texttt{Position +1} next visit).
    \item Never sit with your back to the sea—it implies you do not fear the drowned.
    \item A shared cup of sake before a duel means first blood, not death.
    \item After a battle, offer a handful of salt to the wind before tending your wounds.
\end{itemize}
\textbf{Strings:} Storm‑charm (one weather bypass), broken arrow (proof of oath), ancestor shard (a piece of the drowned kingdom’s crystal), seawater bead (one reroll on a swimming or endurance test).
\end{infobox}

\section{Lethai‑ar — River‑Bound Oath‑Keepers of Ayokha}

\noindent \textit{We do not bargain. We witness. And what we witness, we may choose to remember.}

The Lethai‑ar of Ayokha are the living bridges between mortals and the river spirits that launder the world’s promises. Their ancestors were river‑folk that ascended to the realm of the fair folk, and they now walk the middle path, a step between bodies of water and the world of men. They keep lanterns that never go out, record every contract spoken on the monsoon wharves, and bind their vows with silken cords dyed in river mud.

Unlike their northern kin who swore to distant deities, the Lethai‑ar serve \textbf{Benzaiten} (the river of words) or \textbf{Inari} (the fox’s ledger). Their vows are physical cords, braided from their own hair and dyed with river mud. When a vow is fulfilled, the cord is burned. When it is broken, the cord tightens around the oath‑breaker’s wrist until blood is drawn.

\vspace{0.5em}
\dusana{I once saw a Lethai‑ar witness a contract between a Sihai silk merchant and a Nihori wanderer. The merchant lied about the thread count. The Lethai‑ar said nothing—just cut the witness cord. The merchant’s next shipment arrived as silkworms, not bolts. The wanderer paid double. The Lethai‑ar kept the cord as a reminder.}

\letter{“A witness cord is not a leash. It is a mirror. If you speak truly, you see yourself. If you lie, you see what you could have been. I have cut two cords in my life. I still dream of the faces in the silk.”}{— Lady Suyin, retired Vigil, now keeper of a way‑station for lost travellers}

\begin{infobox}
\textbf{Cultural Mechanics}
\begin{itemize}
    \item \textbf{Witness Clause:} Once per session, when you are present for an oath or contract, you may declare yourself a witness. That oath gains +1 DV to break, and if broken, the oath‑breaker owes you a personal favor (in addition to any other penalty). The witness cord glows faintly for the duration.
    \item \textbf{Lantern of Attendance:} You carry a small oil lamp. As long as it burns, you cannot be surprised in social situations (the flame flickers when deception begins). If the lamp goes out, you lose this benefit for the scene. Refilling the lamp requires a minor offering (a prayer to Benzaiten, a coin in a river).
    \item \textbf{Silk Debt:} Once per arc, you may offer a witness cord as a String to a faction or NPC. While they hold it, you gain +1 die to all social rolls against them, but they may also call on you to witness one of their own oaths (refusing costs the String and 1 Fatigue).
\end{itemize}
\textbf{Etiquette Hooks}
\begin{itemize}
    \item Touch the lamp before speaking – \texttt{Position +1} for truth‑telling.
    \item Never pour the ninth cup (the Vigil enforces the Ayokhan taboo).
    \item A gift of silk is a promise; a gift of silk \textit{cut} is a declaration of feud.
    \item When entering a negotiation, state your name and your witness cord’s color (white for mercy, red for binding, black for truth).
\end{itemize}
\textbf{Strings:} Vigil lantern (social edge), silk witness cord (proof of oath), monsoon mirror (reflects hidden intentions), river stone (one question about a river’s witness).
\end{infobox}

\section{Aelinnel — The Bell‑Tenders of the Inland Sea}

\noindent \textit{We live among the spirits. They borrow milk. We borrow mathematics. It is a fair trade.}

The Aelinnel of the East do not hide from the spirits—they \emph{coexist}. Small spirits inhabit every rock, tree, and wave. The Aelinnel have learned to negotiate with them using copper bells, precise counting, and occasional offerings of pickled plums. They are the engineers of shrines, the builders of floating bridges, and the keepers of the bell‑lines that keep tide‑spirits from straying. Some say Kuva, the sky, taught them the first counting song, and that the Hollow is simply the space where the count fails.

Unlike their western kin, who lean toward the \textbf{Gift of the Mind}, the Eastern Aelinnel practice a folk mathematics that blends logic and superstition. They believe every spirit has a \textit{number‑name}—a frequency that can be rung out of a bell, a sequence of steps that can be walked, a knot that can be tied. To offend a spirit is to miscount.

\vspace{0.5em}
\dusana{“I asked an Aelinnel bell‑tender how she kept the tide‑spirits from flooding her village. She tapped her temple. ‘I count them,’ she said. ‘One, two, three, four, five, six, seven, eight. Then I stop. They know that if I ever count nine, I am not counting them.’”}

\letter{“The Hollow is not a place. It is a miscount. The ninth step that is never there, the ninth cup that is never poured, the ninth name that is never spoken. We do not fight the Hollow. We leave one empty chair, one unlit lamp, one uncounted stone. That is the price of sitting at the table.”}{— Tumble, bell‑warden of the Seventh Tide, from a letter sealed with copper wax}

\begin{infobox}
\textbf{Cultural Mechanics}
\begin{itemize}
    \item \textbf{Bell‑Tending:} Once per scene, when near a shrine or threshold, you may ring a copper bell. The next environmental hazard (tide, fog, landslide) is delayed by 1 segment of the relevant clock. The bell’s tone also pacifies minor spirits (Cap 1) for the scene.
    \item \textbf{Fae Courtesy:} An Aelinnel never says “no” to a spirit. They say “that road is closed for repairs” or “the season is not right.” This phrasing cancels the first Story Beat generated by a spirit encounter. If forced to refuse outright, they mark 1 Fatigue (the spirit’s disappointment is heavy).
    \item \textbf{Kuva’s Count:} Once per journey, you may roll Wits + Lore (DV 3) to recall a forgotten sky‑omen; on success, the next travel leg ignores the first Club card drawn. On a miss, the omen is true, but the GM draws an additional Club for that leg.
    \item \textbf{The Ninth Step:} When you cross a threshold that has a missing ninth step (a stair, a bridge, a door), you may \textit{count} the absence. Gain +1 die to resist the Hollow’s attention for the scene, but mark 1 Obligation (the Hollow now knows your name).
\end{itemize}
\textbf{Etiquette Hooks}
\begin{itemize}
    \item Count your steps on a bridge—odd numbers invite mischief.
    \item Never refuse a spirit’s offered food—but you need not eat it. Holding it is enough.
    \item A copper coin in a shrine’s offering box buys one question about local conditions.
    \item At dusk, ring a bell once, then listen. If you hear a second ring, leave an offering.
\end{itemize}
\textbf{Strings:} Bell‑charter (right to ring a named shrine), tide‑bead (safe passage on Inland Sea), offering‑box key (access to shrine supplies), counting stone (one reroll on a navigation test in fog).
\end{infobox}

\section{Aelaerem — The Hollow’s Neighbors}

\noindent \textit{We live next to the thing that watches from under the hill. We leave butter. It leaves us alone. Mostly.}

The Aelaerem of Nihori and Ayokha have adapted their hedge‑magic to the eastern thresholds. Here the Hollow wears a different face—it is the space between bamboo stalks, the silence after the storm bell, the missing step on the temple stair. They survive through ritual appeasement: a red thread tied at crossroads, a bowl of rice left at dusk, a single bell rung and then muffled. They whisper that the Hollow is the shadow of Kuva, cast when the sky first turned from the sun.

Unlike western Aelaerem, who live in close‑knit villages, the Eastern Aelaerem are often solitary or live in small family groups. They are the ones who know where the boundary between the living and the Hollow is thinnest, and they are the ones paid to keep it closed. Their magic is humble: a knot to bind a promise, a thread to guide a lost ghost, a cup of sake to buy a night’s silence.

\vspace{0.5em}
\dusana{“A halfling once woke me before dawn. ‘The Hollow walked last night,’ she said. I asked how she knew. She pointed to a single red thread tied around my wrist. ‘It moved,’ she said. ‘Sleep with that knot, and it will know where you are.’”}

\letter{“My grandmother taught me that the Hollow is not evil. It is hungry. There is a difference. A hungry thing can be fed. Pour a cup, tie a thread, leave a bell on the doorstep. Do not feed it too much, or it will learn your name.”}{— Old Mila, hedge‑keeper of the Moss Road, from a letter found under a floorboard}

\begin{infobox}
\textbf{Cultural Mechanics}
\begin{itemize}
    \item \textbf{The Empty Cup:} Once per journey, when you are lost or pursued, you may pour a cup of sake or water onto the ground and whisper a name that is not yours. The Hollow will lead you away from danger—but you will forget one minor memory (the GM chooses a trivial fact; if you have no minor memories, advance the Hollow Attention clock [4]). The forgotten memory can be recovered by returning to the same spot and speaking your own name.
    \item \textbf{Neighbor’s Courtesy:} When camping, if you leave a small offering (rice, salt, a button) at the edge of the light, the night’s travel hazard roll is made with Advantage (roll twice, take the less severe). If you forget three nights in a row, the Hollow begins to watch. Start a Hollow Attention clock [4].
    \item \textbf{Kuva’s Eavesdrop:} Once per night, if you sleep under open sky, you may ask one yes/no question about the next day’s weather; the GM answers truthfully, but the question must be whispered to the wind. On a 1, the wind answers but the Hollow also hears.
    \item \textbf{Red Thread Binding:} You may tie a red thread around an object or a threshold to ward against minor spirits (Cap 1) for one scene. If the thread is cut, the spirits notice—start a Spirit Attention clock [4].
\end{itemize}
\textbf{Etiquette Hooks}
\begin{itemize}
    \item Never whistle after dusk—the Hollow thinks you are calling it.
    \item A red thread tied around a lantern means “this light is spoken for.”
    \item Count to eight before opening a door that has been closed for more than a season.
    \item If you meet a stranger at a crossroads at midnight, offer them a cup of sake. If they refuse, walk away without looking back.
\end{itemize}
\textbf{Strings:} Red thread bundle (one safe crossing), Hollow cup (proof of a debt paid in memory), offering plate (right to camp in a dangerous spot), wind‑knot (one question about the Hollow’s mood).
\end{infobox}

\chapter{Travel Decks for the Eastern Realms}

Use the standard Fate’s Edge travel procedure (draw one card from each suit: Spade=place, Heart=actor, Club=pressure, Diamond=leverage). Highest rank sets clock size (4/6/8/10). Each deck below replaces the regional decks from the core travel guide.

\section{Sihai — The Central Kingdom}
\begin{longtable}{c X}
\toprule
\textbf{Spades – Places} & \\
\midrule
2–5   & Examination hall; ink‑stained desks, silent proctors. \\
6–10 & Grand Canal lock; barges queue, officials stamp manifests. \\
J–K  & Ancestral temple; spirit tablets, a single gong. \\
A    & The Forbidden Archive; scrolls that rewrite themselves. \\
\midrule
\textbf{Hearts – People} & \\
\midrule
2–5   & Provincial censor with a list of forbidden words. \\
6–10 & Salt merchant; her scales are honest—except when they aren’t. \\
J–K  & Lethai‑thora memory‑keeper; carries a silence‑token. \\
A    & The Emperor’s Quiet Chancellor; speaks only in couplets. \\
\midrule
\textbf{Clubs – Complications} & \\
\midrule
2–5   & A scroll is found to be a forgery; the party is blamed. \\
6–10 & Monsoon breaks early; the canal closes (delay 2 segments). \\
J–K  & Ancestor demands a public accounting of a forgotten debt. \\
A    & The Mandate wavers; local officials arrest anyone with foreign papers. \\
\midrule
\textbf{Diamonds – Leverage} & \\
\midrule
2–5   & A stamped pass from a provincial governor (one use). \\
6–10 & A copy of the Celestial Ledger (reveals one hidden tax). \\
J–K  & A Lethai‑thora silence‑token (one question the keeper must answer). \\
A    & A junta seal (impersonate an official for one scene). \\
\bottomrule
\end{longtable}

\section{Nihori — The Isles of the Dawn Spirit}
\begin{longtable}{c X}
\toprule
\textbf{Spades – Places} & \\
\midrule
2–5   & A swordsmith’s forge; waterfall cooling tank, prayer strips on the bellows. \\
6–10 & A bamboo grove where duels are fought at dawn. \\
J–K  & A storm‑shrine on a cliff; ropes for climbing during high tide. \\
A    & The Drowned Palace; half‑submerged columns, silent fish. \\
\midrule
\textbf{Hearts – People} & \\
\midrule
2–5   & Wanderer with a broken sword and a grudge. \\
6–10 & Shrine maiden; her bells keep the tide‑spirits at bay. \\
J–K  & Aelinnel bell‑tender; sells tide‑beads. \\
A    & The Storm Emperor’s ghost; offers a deal to the desperate. \\
\midrule
\textbf{Clubs – Complications} & \\
\midrule
2–5   & A volcano god stirs; ashfall covers the path (all actions Desperate). \\
6–10 & A pirate flotilla claims the strait; pay toll or fight. \\
J–K  & The Hollow watches from the bamboo; a character hears their own name whispered. \\
A    & A duel challenge is issued; refusal loses face for the party. \\
\midrule
\textbf{Diamonds – Leverage} & \\
\midrule
2–5   & A swordsmith’s token (one reroll on a failed melee attack). \\
6–10 & A shrine bell that reveals hidden spirits for one scene. \\
J–K  & A wanderer’s oath‑bracelet (one use to compel a service). \\
A    & A storm‑charmer’s hat (safe passage through one storm). \\
\bottomrule
\end{longtable}

\section{Ayokha — The Monsoon Throne}
\begin{longtable}{c X}
\toprule
\textbf{Spades – Places} & \\
\midrule
2–5   & A river market on stilts; goods in baskets, prices in smiles. \\
6–10 & A step‑pyramid temple; carvings of the monsoon’s first year. \\
J–K  & A silk village; mulberry groves, spinning sheds. \\
A    & The Bridge of Three Echoes; appears only once a century. \\
\midrule
\textbf{Hearts – People} & \\
\midrule
2–5   & River merchant whose boat has more secrets than cargo. \\
6–10 & A spirit‑medium who speaks for the \textit{nat} of the ford. \\
J–K  & Lethai‑ar silk‑weaver; offers a witness clause. \\
A    & The God‑King’s astrologer; reads omens in the monsoon’s arrival. \\
\midrule
\textbf{Clubs – Complications} & \\
\midrule
2–5   & The river floods; a bridge is swept away. \\
6–10 & A rival merchant spreads rumors of cursed silk. \\
J–K  & The \textit{nat} of the well is offended; no one may draw water until appeased. \\
A    & The God‑King falls ill; all travel grinds to a halt. \\
\midrule
\textbf{Diamonds – Leverage} & \\
\midrule
2–5   & A monsoon schedule (predicts safe sailing days). \\
6–10 & A silk voucher (trade goods at 10\% profit). \\
J–K  & Lethai‑ar witness cord (enforce any oath). \\
A    & A blessing from the God‑King (one national favor). \\
\bottomrule
\end{longtable}

\section{The Inland Sea (Shared Deck)}
\begin{longtable}{c X}
\toprule
\textbf{Spades – Places} & \\
\midrule
2–5   & A floating shrine; bells on every corner. \\
6–10 & A pirate cove disguised as a fishing village. \\
J–K  & A reef where spirit‑ships are said to anchor. \\
A    & The sunken gate of the drowned Lethai kingdom. \\
\midrule
\textbf{Hearts – People} & \\
\midrule
2–5   & A pilot with a cloud‑chart and a price. \\
6–10 & A Kuvani horse‑trader bound for Sihai. \\
J–K  & A Tulkani storyteller with a vow‑cord; she whispers a sky‑omen. \\
A    & The Drowned Captain; offers passage for a memory. \\
\midrule
\textbf{Clubs – Complications} & \\
\midrule
2–5   & A sudden squall; all navigation rolls at disadvantage. \\
6–10 & A pirate attack; pay cargo or fight. \\
J–K  & A spirit‑whirlpool; lose 1 Supply to escape. \\
A    & The Drowned Kingdom rises; all spirits of the deep demand a tithe. \\
\midrule
\textbf{Diamonds – Leverage} & \\
\midrule
2–5   & A pilot’s chart (one leg ignores weather). \\
6–10 & A pirate’s flag (safe passage through one blockade). \\
J–K  & A Tulkani story‑knot (one reroll on a failed negotiation). \\
A    & A memory of the drowned (ask one historical question; answer is true). \\
\bottomrule
\end{longtable}

\chapter{Patrons of the Eastern Realms}

\section{Shennong — The Farmer of a Thousand Grains}
\noindent\textbf{Domain:} Agriculture, medicine, the rhythm of planting and harvest.\\
\textbf{Lore:} The spirit who taught the first emperor to plow. He is patient, generous, and furious about waste. His shrines are terraced fields, not temples. His offerings are the first handful of grain and the last. He is said to have tasted a hundred herbs to discover which heal and which poison; his followers still test medicines on themselves before giving them to others.

\vspace{0.5em}
\dusana{A priest of Shennong once refused to heal a wound because the injured man had burned a field. “The grain remembers,” he said. “So do I.”}

\letter{“Shennong does not speak in words. He speaks in the taste of soil and the colour of the sprout. A healthy rice stalk is a sentence. A blighted one is a warning. Learn to read the field, and you will never need a temple.”}{— Yi, terraced farmer of the upper Sihon, from a letter stained with mud}

\begin{infobox}
\textbf{Rites}
\begin{itemize}
    \item \textit{Low – The Seed’s Memory (4 XP):} Touch a grain; know where it was grown and what season. (Mark +1 Obligation) The grain’s memory may include the farmer’s name, the weather, or a debt unpaid.
    \item \textit{Standard – Shennong’s Brew (8 XP):} Create a medicinal tea that heals 1 Harm or removes one Poison. The tea tastes of earth and requires a pinch of the patient’s own sweat. (Mark +2 Obligation)
    \item \textit{High – The Thriving Field (14 XP):} Double the yield of one farm for a season; the land is tired for the next year. Shennong demands a promise that the surplus will feed the hungry, not be hoarded. (Mark +3 Obligation; if you break the promise, the land fails the next season.)
\end{itemize}
\end{infobox}

\section{Amaterasu — The Hidden Sun}
\noindent\textbf{Domain:} Light, truth, the sun behind clouds.\\
\textbf{Lore:} The Lethai‑thora speak of a sun that once hid in a cave. Her return is a promise that truth cannot be hidden forever. She is served by mirror‑keepers and lamp‑lighters. Her shrines are dark, windowless rooms with a single bronze mirror; the worshipper must light a candle and watch until they see their own face clearly. That moment is her blessing.

\textbf{Favored Deeds:} Revealing a lie, lighting a lamp in a dark place, honoring a promise made at dawn, polishing a mirror.\\
\textbf{Taboos:} Hiding a crime, extinguishing a sacred flame, swearing an oath you do not mean, breaking a mirror without ritual.

\letter{“I polished the mirror for forty days before I saw her. She looked nothing like the paintings. Her face was my own, but younger, and she was laughing at something I had forgotten. That is her blessing: not truth, but the memory of truth.”}{— Hoshi, lamp‑lighter of the Western Terrace, from a letter found in a shrine offering box}

\begin{infobox}
\textbf{Rites}
\begin{itemize}
    \item \textit{Low – The Crack of Light (4 XP):} A single beam of sunlight illuminates a shadowed area for one scene. It lasts only as long as you hold a polished metal surface toward the sun. (Mark +1 Fatigue)
    \item \textit{Standard – The Unhidden Ledger (8 XP):} One hidden document or secret becomes visible to you for a moment. You must speak the secret aloud to the first person you meet, or the vision scars your eyes (–1 die to Perception for a day). (Mark +2 Obligation)
    \item \textit{High – The Sun’s Return (14 XP):} Banish one supernatural darkness or illusion permanently. The ritual requires a mirror, dawn light, and the true name of the darkness. (Mark +3 Corruption; your shadow is slightly shorter for a season.)
\end{itemize}
\end{infobox}

\section{Susanoo — The Storm That Clears}
\noindent\textbf{Domain:} Storms, courage, the violence that resets the board.\\
\textbf{Lore:} A troubled spirit who destroys to make way for the new. His followers are ronin, exiles, and those who have lost everything. He is not evil—only furious. He respects those who face him without flinching. His shrines are hilltops struck by lightning, and his offerings are broken weapons and torn banners.

\textbf{Favored Deeds:} Challenging a tyrant, surviving a storm without shelter, cutting a knot that binds unjustly, repairing a broken bridge after crossing it.\\
\textbf{Taboos:} Attacking the helpless, breaking a weapon in anger, fleeing a duel, sheltering from a storm when others are outside.

\letter{“My father prayed to Susanoo before every duel. He never won. But he never ran. That is the god’s only demand: that when the lightning comes, you stand where it can see you.”}{— Sato, ronin of no fixed name, from a letter pinned to a shrine gate}

\begin{infobox}
\textbf{Rites}
\begin{itemize}
    \item \textit{Low – The Squall’s Edge (4 XP):} Your next melee attack gains +1 Harm, but you cannot retreat this scene. The attack leaves a faint scorch mark shaped like a lightning bolt. (Mark +1 Fatigue)
    \item \textit{Standard – Susanoo’s Breath (8 XP):} Summon a gust that extinguishes all non‑magical fires in a zone. The wind smells of sea salt and wet pine. (Mark +2 Obligation)
    \item \textit{High – The Clearing Storm (14 XP):} End one ongoing conflict (battle, chase, argument) by force—both sides take Harm 1. The violence is not merciful, but it is final. (Mark +3 Corruption; your voice grows hoarse for a week.)
\end{itemize}
\end{infobox}

\section{Benzaiten — The River of Words}
\noindent\textbf{Domain:} Speech, music, the flow of negotiation.\\
\textbf{Lore:} A spirit of the river’s sound. She favors poets, merchants, and those who can speak three languages. Her shrines are on bridges, where the water’s voice carries the words of travellers. She is said to have taught the first courtesan the art of saying “yes” and meaning “no.” Her followers are often diplomats, spies, and lovers.

\textbf{Favored Deeds:} Translating a misunderstood agreement, singing a lullaby to a spirit, resolving a dispute without bloodshed, writing a poem that makes a stranger weep.\\
\textbf{Taboos:} Interrupting a poem, silencing a witness, lying about the meaning of a word, using a song to harm.

\letter{“Benzaiten came to me in the sound of water over stones. She said: ‘Every word is a current. Speak gently, and the river will carry you. Speak cruelly, and you will drown.’ I have never lied since. I have also never been silent.”}{— Paya, silk trader of Ayokha, from a letter to her daughter}

\begin{infobox}
\textbf{Rites}
\begin{itemize}
    \item \textit{Low – The Flowing Syllable (4 XP):} Gain +1 die to Sway for one scene; your voice carries unnaturally far. The words leave a faint silver echo that fades after a minute. (Mark +1 Fatigue)
    \item \textit{Standard – The Shared Tongue (8 XP):} You and one target understand each other perfectly for one scene, regardless of language. If either of you lies, the other hears a discordant note—like a bell cracking. (Mark +2 Obligation)
    \item \textit{High – The River’s Verdict (14 XP):} Compel two parties to accept a binding arbitration; refusal causes –2 dice to all social rolls for a week. The verdict must be spoken while both parties touch the same body of water. (Mark +3 Obligation)
\end{itemize}
\end{infobox}

\section{Inari — The Fox’s Ledger}
\noindent\textbf{Domain:} Commerce, secrets, the threshold between realms.\\
\textbf{Lore:} A spirit of the marketplace who wears a fox mask. She is generous to those who honor her rites and merciless to those who cheat. Her shrines are tucked between shops, often with a small bronze fox at the door. She does not speak; she writes on ledgers that appear overnight.

\textbf{Favored Deeds:} Paying a fair price, sharing a meal with a stranger, keeping a promise made over a sale, refusing to sell something that would cause harm.\\
\textbf{Taboos:} Selling a known counterfeit, hoarding a surplus, refusing to bargain, lying about a product’s origin.

\letter{“Inari left a ledger on my pillow. It listed every customer who had ever shortchanged me, and beside each name, a price: a lost coin, a broken cart, a spoiled harvest. I did not collect. I burned the ledger. The next morning, a fox lay dead on my doorstep. I understood: Inari does not offer justice. She offers the knife. You must choose whether to use it.”}{— Merchant Ken, from a letter found in his abandoned shop}

\begin{infobox}
\textbf{Rites}
\begin{itemize}
    \item \textit{Low – The Mask’s Glimpse (4 XP):} For one scene, you can see if a person is lying about money or trade. Their ears glow faintly red when they lie. (Mark +1 Obligation)
    \item \textit{Standard – Inari’s Scale (8 XP):} Guarantee a 10 % profit on a trade; on a miss, the profit turns to sand. The sand is warm and smells of cinnamon. (Mark +2 Obligation)
    \item \textit{High – The Fox’s Market (14 XP):} For one hour, any mundane item can be bought or sold in any settlement, regardless of availability. The prices are strange—a memory for a sword, a year of silence for a horse. (Mark +3 Obligation; Inari will collect her due within the season.)
\end{itemize}
\end{infobox}

\section{Hachiman — The Bow of Ancestors}
\noindent\textbf{Domain:} War, archery, protection of the homeland.\\
\textbf{Lore:} A warrior spirit who defends borders. His followers are mountain guards, frontier ronin, and those who have lost family to invasion. He is said to ride a black horse and carry a bow that whispers the names of fallen soldiers. His shrines are stone towers at mountain passes, with a single arrow embedded in the door.

\textbf{Favored Deeds:} Defending a village, training a student, using a bow only when necessary, returning a captured banner to its people.\\
\textbf{Taboos:} Shooting a civilian, abandoning a post, breaking a bow in anger, killing a supplicant who asks for mercy.

\letter{“Hachiman does not answer prayers for victory. He answers prayers for courage. I asked him to let me live through the battle. He showed me a vision of my son growing up without a father. I chose to live anyway. The arrow passed through my shoulder, not my heart. That is his mercy.”}{— Commander Toshiro, from a letter to his wife, never sent}

\begin{infobox}
\textbf{Rites}
\begin{itemize}
    \item \textit{Low – The Arrow’s Memory (4 XP):} You know the last person your weapon wounded. The knowledge comes as a brief vision of their face. (Mark +1 Fatigue)
    \item \textit{Standard – Hachiman’s Wall (8 XP):} For one scene, you and adjacent allies gain +1 Armor against ranged attacks. The arrows that would have struck you fall short, as if stopped by an invisible wind. (Mark +2 Obligation)
    \item \textit{High – The Ancestor’s Shot (14 XP):} One arrow ignores all mundane cover and armor. The arrow leaves a faint blue light in its wake, visible to anyone who has lost a family member to war. (Mark +3 Obligation; the arrow cannot be retrieved.)
\end{itemize}
\end{infobox}

\section{Kuva — The Sky That Takes Many Names}
\noindent\textbf{Domain:} The sky, storms, stars, open air, witness of oaths.\\
\textbf{Lore:} Kuva is the dome above all roads—the endless blue, the monsoon wrath, the night’s silent gaze. To the Kuvani of the eastern steppe, Kuva is a stern witness who hears every promise made under open heaven. To the Dhaharan of the fertile lowlands, Kuva is a storm‑mother, destroyer and renewer, whose rains demand sacrifice and give life. To the Tulkani wanderers, Kuva is a shape‑changing presence—neither kind nor cruel, only watchful, and a slighted sky brings ill fortune. To the Aelinnel, Kuva is the first counter, the one who taught the counting song. To the Aelaerem, Kuva is the Hollow’s boundary—the sky is where the Hollow’s gaze stops.

\vspace{0.5em}
\dusana{“The Kuvani pour milk upward for Kuva. The Dhaharan offer blood at monsoon bells. The Tulkani tie a knot, throw ash, and look up. The Aelinnel count clouds. The Aelaerem whisper to the wind. One sky, five peoples. Kuva does not correct them. Kuva watches.”}

\letter{“Kuva has no shrine because Kuva is the roof of every shrine. When I pray, I do not bow. I look up. The sky does not answer in words. It answers in weather. A clear day means ‘go.’ A storm means ‘wait.’ Lightning means ‘run.’ Learn to read the weather, and you will never need a priest.”}{— Rani, Kuvani herder, from a letter scratched on a tent flap}

\letter{“We Dhaharan say that Kuva’s third eye is the monsoon. When it opens, the floods come. When it closes, the drought begins. We do not pray for her to close or open. We pray for her to blink.”}{— The monk at the Temple of the Third Rain, from a fragment of palm leaf}

\begin{infobox}
\textbf{Rites}
\begin{itemize}
    \item \textit{Low – The Sky’s Witness (4 XP):} For one scene, when you speak a promise outdoors, any listener who hears it must test Resolve (DV 3) to lie about that promise. (Mark +1 Obligation; Kuva remembers your own words as well.) If you break your own promise, the GM gains 2 Story Beats to spend on a sky‑related complication (lightning, sudden fog, a hawk stealing your food).
    \item \textit{Standard – Kuva’s Gaze (8 XP):} Call down a brief flash of lightning or a clap of thunder. All who see or hear it must test Spirit (DV 4) or be Stunned for one exchange. (Mark +2 Obligation; cannot be used indoors or underground.) The lightning always strikes the highest point in the zone, which may be inconvenient.
    \item \textit{High – The Monsoon’s Due (14 XP):} Once per season, you may pray for rain (or clear skies) over a one‑mile area. The weather shifts within an hour and lasts for a day. The cost is a visible scar on your own flesh—a mark that Kuva has accepted your debt. (Mark +3 Obligation; the scar may be noticed by others who know Kuva’s signs.) If you pray while standing under open sky, you may also ask for a specific omen—a double rainbow, a lightning strike on a target, a wind that carries a message.
\end{itemize}
\end{infobox}

\noindent\textit{Cultural Variations}
\begin{itemize}
    \item \textbf{Kuvani:} Kuva is seen as a watcher of oaths. Offerings are mare’s milk thrown upward or a knot tied in a wind‑charm. The debt is \emph{visibility}: you are accountable because you are seen.
    \item \textbf{Dhaharan:} Kuva is a storm‑mother of many arms, wearer of a third eye. Offerings are blood (a chicken, or rarely stronger). The debt is \emph{cyclical sacrifice}: pay what you owe, or Kuva takes from you—crops, memory, or kin.
    \item \textbf{Tulkani:} Kuva shifts between faces. Offerings are ash thrown upward or a breath spoken into a knot. Kuva’s first cast shadow is called \emph{Ikasha}—the original darkness formed when the sky turned away from the sun. The debt is \emph{awareness}: simply looking up and acknowledging the sky is a form of payment.
    \item \textbf{Aelinnel:} Kuva is the Great Accountant. Offerings are a counted sequence (eight beads, eight stones, eight breaths). The debt is \emph{miscount}: if you ever count to nine where Kuva can hear, the sky will answer.
    \item \textbf{Aelaerem:} Kuva is the boundary. Offerings are silence. The debt is \emph{attention}: if you look up and Kuva looks back, the Hollow remembers you.
\end{itemize}

\chapter{Adventure Seeds}
\begin{enumerate}
    \item \textbf{The Forgotten Mandate (Sihai):} A Lethai‑thora memory‑keeper is assassinated. The party must retrieve her memory‑silence before the Censorate burns it. The secret: an emperor faked his own death.
    \item \textbf{The Wanderer’s Grave (Nihori):} A duel was fought to a draw. Now both families demand the party arbitrate the next match—while the ghost of the first duel watches.
    \item \textbf{The Monsoon Seal (Ayokha):} A merchant’s seal is stolen the day before the first monsoon ships arrive. The party has one day to find it—or the merchant’s family will be sold.
    \item \textbf{The Bell That Did Not Ring (Inland Sea):} A shrine’s bell falls silent. Without its chime, the tide‑spirits are confused. Two villages prepare to blame each other.
    \item \textbf{The Fox’s Bargain (Sihai):} A farmer claims Inari herself visited him. She promised a great harvest—in exchange for his daughter’s name. The daughter has forgotten who she is.
    \item \textbf{The Storm Emperor’s Armor (Nihori):} A set of legendary armor is discovered in a sunken tomb. Three clans claim it. The party must authenticate—or forge—provenance.
    \item \textbf{The Silk Witness (Ayokha):} A Lethai‑ar witness cord is found cut. Without it, a trade treaty cannot be enforced. The party must re‑weave the cord—and the truth.
    \item \textbf{The Hollow Well (Mixed):} A well in a border village is dry. The Aelaerem say the Hollow drank it. The party must descend to negotiate.
    \item \textbf{The Copper Bell (Aelinnel):} A shrine’s bell‑tuning has gone wrong. The spirits are agitated. The party must re‑count the steps—or appease them with a story.
    \item \textbf{The Ninth Step (Any):} A temple stair has a missing ninth step. Travelers who count “eight” and stop are safe. Those who try to find the ninth invite the Hollow. The party must enter the Hollow to retrieve a lost child.
\end{enumerate}

\chapter{GM Toolkit: Eastern Clocks and Fronts}
\begin{longtable}{X c c X}
\toprule
\textbf{Clock} & \textbf{Region} & \textbf{Segments} & \textbf{On Fill} \\
\midrule
Mandate Frays & Sihai & 6 & Civil war; passes close \\
Storm Season & Nihori & 8 & Typhoon; all travel Desperate for 2 legs \\
Monsoon Delay & Ayokha & 6 & Famine; grain prices double \\
Hollow Stirring & Any with Aelaerem & 4 & Something follows the party \\
Sky Debt & Any with Kuva & 6 & A weather catastrophe (drought, flood, lightning plague) \\
\bottomrule
\end{longtable}

\subsection*{Faction Influence (d6 per region)}
\textbf{Sihai:}
\begin{enumerate}[label=\arabic*.]
    \item Censorate (hunts truth)
    \item Salt Monopoly (controls trade)
    \item Lethai‑thora memory‑keepers (hidden knowledge)
    \item Provincial Warlords (armed neutrality)
    \item Emperor’s Court (nominal authority)
    \item Aelinnel bell‑tenders (infrastructure)
\end{enumerate}

\textbf{Nihori:}
\begin{enumerate}[label=\arabic*.]
    \item Ashigaru League (commoner levies)
    \item Wanderer Brotherhood (masterless swords)
    \item Shrine Maidens (spirit diplomacy)
    \item Pirate Clans (coastal pressure)
    \item Storm Emperor Revivalists (lost dynasty)
    \item Aelaerem Hollow‑Watchers (threshold guardians)
\end{enumerate}

\textbf{Ayokha:}
\begin{enumerate}[label=\arabic*.]
    \item God‑King’s Astrologers (ritual authority)
    \item River Merchant Guild (trade power)
    \item Lethai‑ar Silent Vigil (oath enforcement)
    \item Spirit Mediums (\textit{nat} negotiation)
    \item Monsoon Captains (naval strength)
    \item Mixed‑Heritage Freedmen (cultural bridge)
\end{enumerate}

\chapter{Safety and Consent at the Table}
This expansion deals with themes of spirit possession, ancestor veneration, and cultural practices inspired by East and Southeast Asian traditions. These elements are intended to create atmosphere and tension, not to exoticise or stereotype.

Before play, establish:
\begin{itemize}
    \item \textbf{Lines:} Hard boundaries. Example: no forced spirit possession of PCs, no harm to shrine animals, no desecration of ancestral tablets.
    \item \textbf{Veils:} Soft boundaries. Example: possession can be described obliquely (“the spirit speaks through them”) rather than depicted as loss of agency.
    \item \textbf{X‑Card:} Any player may tap a card marked X to pause or remove a scene without explanation. The table pivots immediately.
\end{itemize}

Use these tools. Trust them.

\chapter{Colophon}
This expansion was written by the edge of the Inland Sea, on a page that refused to dry. The Aelinnel bell‑ringer said I had counted to nine. I apologized three times. He accepted the first apology, ignored the second, and used the third to tune his next bell.

\vspace{1em}
\dusana{At the floating shrine of the Seventh Tide, the lamp is borrowed, the ink is squid, and the Hollow is not hungry—today.}

\end{document}
```